\documentclass[12pt,letterpaper]{article}
\usepackage[utf8]{inputenc}
\usepackage[spanish]{babel}
\usepackage{amsmath}
\usepackage{amsfonts}
\usepackage{amssymb}
\usepackage[margin=2.5cm]{geometry}

\begin{document}

\begin{center}
    \Large \textbf{Evaluación de Examen 2: Unidad II} \\
    \large Física para Ciencias de la Salud (005-1713) \\
    \vspace{0.5cm}
    \textbf{Estudiante:} Franyerika Rodríguez \quad \textbf{C.I.:} No indicado
    \vspace{0.3cm} \\
    \large \textbf{Calificación Final: 10/10 puntos}
\end{center}

\vspace{0.5cm}

\section*{Análisis de Resultados}

\subsection*{1. Cinemática (Aceleración media)}
\textbf{Procedimiento y resultado:} Extrae de manera impecable los datos iniciales y la fórmula de aceleración ($a = \frac{v_f - v_i}{t}$). Sustituye los valores correctamente ($0,6 / 0,15 = 4$) y, de manera sobresaliente, realiza un análisis dimensional gráfico a la derecha para demostrar cómo se cancelan las unidades hasta obtener $\text{m/s}^2$. \\
\textbf{Veredicto:} \textbf{Correcto} (2/2 pts).

\subsection*{2. Dinámica (Leyes de Newton)}
\textbf{Procedimiento y resultado:} Plantea la Segunda Ley de Newton de manera vectorial y realiza el despeje algebraico correcto para la aceleración ($a = \frac{F}{m}$). Sustituye los datos numéricos ($150 \text{ N} / 25 \text{ kg}$) para obtener la aceleración correcta de $6 \text{ m/s}^2$. \\
\textbf{Veredicto:} \textbf{Correcto} (2/2 pts).

\subsection*{3. Trabajo Mecánico}
\textbf{Procedimiento y resultado:} Extrae los datos y aplica la fórmula general del trabajo incluyendo el componente trigonométrico ($W = F \cdot d \cdot \cos(\theta)$). Identifica implícitamente que el ángulo es $0^\circ$ e indica correctamente el cálculo final, obteniendo de forma exacta $320 \text{ Joules}$. \\
\textbf{Veredicto:} \textbf{Correcto} (2/2 pts).

\subsection*{4. Energía Potencial}
\textbf{Procedimiento y resultado:} Extrae y lista rigurosamente los datos del problema (masa, altura y gravedad). Aplica correctamente la ecuación de energía potencial gravitatoria ($E_p = m \cdot g \cdot h$), lo cual la lleva al producto matemáticamente exacto de $17,64 \text{ Joules}$. \\
\textbf{Veredicto:} \textbf{Correcto} (2/2 pts).

\subsection*{5. Potencia Mecánica}
\textbf{Procedimiento y resultado:} Plantea directamente la fórmula correcta para la potencia ($P = \frac{W}{t}$) y la relaciona con los datos previamente extraídos. Efectúa la división de $75 \text{ J}$ entre $60 \text{ s}$ obteniendo de manera perfecta el resultado de $1,25 \text{ Watts}$. \\
\textbf{Veredicto:} \textbf{Correcto} (2/2 pts).

\vspace{0.5cm}
\section*{Comentarios Generales y Sugerencias}
¡Felicidades por alcanzar la nota máxima! El examen demuestra un entendimiento sólido tanto de los conceptos físicos como del procedimiento matemático necesario para demostrarlos.
\begin{itemize}
    \item El formato de resolución utilizado, separando el área de trabajo en \textit{Datos} y \textit{Desarrollo}, es óptimo para la organización en ciencias de la salud.
    \item Se valora extremadamente positivo el uso del \textbf{análisis dimensional explícito} (como el efectuado en la pregunta 1 para demostrar el origen del $\text{m/s}^2$); es un excelente hábito analítico.
    \item \textbf{Pequeña sugerencia de forma:} Se recomienda cuidar la legibilidad en el trazado de los resultados finales para evitar confusiones de lectura (como en las unidades finales de la pregunta 3 o las anotaciones de notación científica extras en la pregunta 4), si bien los desarrollos son perfectamente válidos.
    \item ¡Sigue con este fantástico nivel de dedicación!
\end{itemize}

\end{document}